\documentclass[12pt, a4paper]{article}
\usepackage[margin=2.5cm]{geometry}
\usepackage{graphicx}
\usepackage{setspace}
\usepackage{parskip}
\usepackage{titlesec}
\usepackage{hyperref}
\usepackage{apacite}
\usepackage{natbib}
\usepackage{times}
\usepackage{microtype}
\usepackage[none]{hyphenat}
\tolerance=1
\emergencystretch=\maxdimen
\hyphenpenalty=10000
\hbadness=10000

\onehalfspacing

\title{\textbf{Blockmediary: Legal and Compliance Risk Assessment}}
\author{}
\date{BEEM063 FinTech Hackathon \\ University of Exeter}

\begin{document}

\maketitle
\newpage

%---------------------------------------------------------------
\section*{Introduction}
%---------------------------------------------------------------

Blockmediary is a programmable documentary escrow layer for small and medium sized
enterprise cross border trade. The product locks stablecoin, specifically USDC on the
Base network, into a smart contract, verifies the trade documents a seller submits
against agreed release conditions, and releases the funds once those conditions are
satisfied. This mechanism creates regulatory surface area across six distinct domains:
crypto asset and stablecoin regulation, payment services law, anti money laundering and
know your business and know your customer obligations, data privacy law, trade finance
law, and consumer and business protection law. The purpose of this assessment is to work
through each domain in turn, across the jurisdictions Blockmediary has identified as
realistic near term markets, namely the United Kingdom, the European Union, and the
United Arab Emirates, and to record why Turkey has been excluded as a target market
despite the size of its trading relationship with the Gulf region.

The assessment is organised as follows. The first section addresses the regulatory
position in the United Kingdom, the second addresses the European Union, and the third
addresses Turkey briefly, as an excluded jurisdiction rather than an active one. The
fourth section addresses the United Arab Emirates, distinguishing between the federal
mainland regime and the more accessible free zone frameworks. The fifth section addresses
risks that cut across all jurisdictions rather than belonging to any single regulator, such
as sanctions screening, trade based money laundering, and the legal standards that govern
documentary release itself. The sixth section consolidates the regulatory instruments
referenced throughout the assessment into a single index. The seventh section compares
Blockmediary's compliance posture against the closest competing platforms. The eighth and
final section sets out the confirmed United Arab Emirates launch route among the three
principal Emirati jurisdictions available to a fintech entrant, and records why the
earlier choice was superseded. This document is a working compliance reference prepared
for an academic hackathon project and does not constitute legal advice; qualified legal
counsel should be sought before Blockmediary operates in any of the jurisdictions discussed
with real funds or real users.

Blockmediary has validated the core escrow mechanism on testnet as a proof of concept.
Mainnet deployment for real value would additionally require licensing or Virtual Asset
Service Provider registration, resolution of client asset custody status, sanctions
screening, resolution of the General Data Protection Regulation question raised by the
on-chain specification hash, and an independent security audit. These requirements are set
out throughout this assessment as a forward looking compliance roadmap, not as requirements
already satisfied.

%---------------------------------------------------------------
\section*{1. United Kingdom Regulatory Risks}
%---------------------------------------------------------------

\subsection*{1.1 Stablecoin Issuance and Cryptoasset Custody}

The Financial Conduct Authority published consultation paper CP25/14 in May 2025, setting
out proposed rules for firms that issue qualifying stablecoins and for firms that provide
custody of cryptoassets \citep{fcacp2514}. The final rules are expected during 2026 and the
regime will take effect once it clears parliamentary approval. Blockmediary does not issue
USDC, Circle does, so the issuance limb of CP25/14 is not directly engaged. The custody
limb is a closer question. If Blockmediary's smart contract ever holds USDC on behalf of a
buyer or seller in a manner that amounts to safeguarding a cryptoasset as a regulated
activity, the proposed client asset framework for cryptoassets would require authorisation,
segregation of customer assets, and custody arrangements that meet the Financial Conduct
Authority's standards. Blockmediary's response is architectural rather than purely legal.
By using a smart contract escrow model in which funds sit in the on chain contract rather
than in a wallet Blockmediary itself controls, the product moves away from the clearest
form of custody, but this is a technical description rather than a complete regulatory
conclusion. A subsequent readiness review of the United Arab Emirates pathway reached the
same view and it is adopted here for all jurisdictions: regulators generally examine who
can cause, prevent, reverse, upgrade or recover a transfer, and who bears responsibility to
the client, so Blockmediary should not describe itself as definitively non custodial until
that control analysis over the release key, contract administration and recovery rights has
been documented \citep{difcvarareport2026}. For production use, the safer path is to
partner with an authorised cryptoasset custodian for any component of the flow that does
resemble safeguarding. For the hackathon build, the product operates on testnet with mock
USDC, so no regulated custody activity is actually triggered, but the Policy Statement
expected in the summer of 2026 should be monitored and the architecture reassessed once the
final rules are published.

\subsection*{1.2 Crypto Asset Service Provider Licensing}

Consultation paper CP25/40 sets out the Financial Conduct Authority's broader framework
for regulating cryptoasset activities, extending the existing Financial Services and
Markets Act 2000 regime to cover exchange, custody, and related cryptoasset services
provided to United Kingdom persons \citep{fcacp2540}. Whether Blockmediary needs to
register or seek authorisation under this expanded regime depends on how its role in
arranging or operating the escrow is characterised. The existing cryptoasset registration
regime under the Money Laundering Regulations focuses narrowly on anti money laundering
compliance, but the new framework introduces conduct, prudential, and consumer protection
obligations that go considerably further. Blockmediary's approach is to take legal advice
on whether its activity constitutes a regulated cryptoasset service under the new regime,
while operating in testnet only during the hackathon phase so that no United Kingdom
persons transact with real funds. For production, the team should budget for a
twelve to eighteen month authorisation timeline, or alternatively consider launching first
in a free zone jurisdiction such as the Abu Dhabi Global Market or the Dubai International
Financial Centre, where the path to authorisation is comparatively cleaner, before
returning to the United Kingdom market once the domestic regime has settled.

\subsection*{1.3 Payment Services}

The Payment Services Regulations 2017 implement the European Union's second Payment
Services Directive into United Kingdom law and regulate the provision of payment services
\citep{ukpsr2017}. If Blockmediary's escrow function is characterised as holding funds on
behalf of parties, or as executing payment transactions, it could constitute a payment
service requiring authorisation as a payment institution or an electronic money
institution. The smart contract escrow model partially mitigates this concern because
funds pass directly into the contract rather than through an account Blockmediary
operates, but the release instruction, the step at which Blockmediary authorises the on
chain transfer, may still be characterised as execution of a payment transaction. The
mitigation is twofold. First, Blockmediary should obtain a legal opinion on whether the
release instruction itself constitutes execution of a payment transaction under the
regulations. Second, in all product communications and legal documentation, Blockmediary's
role should be framed consistently as verification and orchestration rather than as
payment execution, and for production deployment the fiat on and off ramp layer should be
delivered through a partnership with an already licensed payment institution rather than
by Blockmediary seeking its own licence for that specific function.

\subsection*{1.4 Anti Money Laundering, Know Your Business, and Know Your Customer}

The Money Laundering, Terrorist Financing and Transfer of Funds Regulations 2017, as
amended, require cryptoasset businesses to apply customer due diligence, transaction
monitoring, and suspicious activity reporting \citep{mlr2017}, and the European Union's
sixth Anti Money Laundering Directive tightens the equivalent obligations across the
European Union \citep{amld6}. Underlying this procedural framework is the Proceeds of
Crime Act 2002, which makes money laundering a criminal offence in the United Kingdom and
creates the money laundering reporting officer's mandatory suspicious activity report
obligation under section 330 \citep{poca2002}. A failure to file a suspicious activity
report where one is required is a criminal offence for both the firm and the individual
reporting officer, and the Terrorism Act 2000 creates a parallel disclosure obligation for
suspected terrorist financing \citep{terrorismact2000}. Blockmediary is a virtual asset
service provider under the Financial Action Task Force's definition, since it facilitates
the transfer and custody of virtual assets, and this triggers the full suite of anti money
laundering and counter financing of terrorism obligations. Trade based money laundering is
a particular concern in cross border trade escrow, since over invoicing, under invoicing,
phantom shipments, and document fraud are established typologies within this sector.
Blockmediary's mitigation is to implement a compliance gate before any escrow is funded,
requiring both parties to complete know your business and know your customer checks,
sanctions screening, and wallet screening, using a reputable provider such as Sumsub,
Jumio, or Onfido, and to maintain a trade based money laundering red flag checklist that
watches for invoice prices inconsistent with market rates, unusual shipment routes, vague
goods descriptions, and discrepancies across submitted documents. A United Kingdom money
laundering reporting officer should be appointed for production operation, bearing in mind
that this individual carries personal criminal liability under section 330 of the Proceeds
of Crime Act for failure to disclose. A related complication arises where a third party
platform or freight forwarder initiates a deal on behalf of a buyer and seller rather than
the principals approaching Blockmediary directly. In that scenario, Blockmediary must
ensure that know your customer and know your business checks on the underlying principals
are performed either directly by Blockmediary or by a third party acting under regulation
39 of the Money Laundering Regulations, which permits reliance on third parties subject to
strict conditions: the third party must itself be subject to equivalent anti money
laundering obligations, a written reliance agreement must be in place, and Blockmediary
must be able to obtain the underlying customer due diligence data within two business days
of request. A platform initiated deal should never be permitted to fund before the
underlying buyer and seller have been verified under one of these two routes.

\subsection*{1.5 The Financial Action Task Force Travel Rule}

Recommendation 16 of the Financial Action Task Force, as revised in June 2025, requires
virtual asset service providers to collect, verify, and transmit originator and
beneficiary information alongside virtual asset transfers above a threshold that in most
jurisdictions sits at the equivalent of one thousand United States dollars or euros
\citep{fatf2025}. Every escrow funding and release transaction that Blockmediary processes
is a virtual asset transfer for these purposes, so the Travel Rule requires Blockmediary to
collect and, where the counterparty is itself a virtual asset service provider, share
identifying information about both parties. The revised standards take full effect
domestically across most jurisdictions by the end of 2030, but the United Kingdom, the
European Union under the Markets in Crypto Assets Regulation, and the United Arab Emirates
are already implementing equivalent requirements ahead of that deadline. Because
Blockmediary already collects full know your business and know your customer data at
onboarding as part of its anti money laundering programme, this same dataset satisfies the
Travel Rule's data requirement, and for wallet to contract transfers the product should
implement Travel Rule data transmission to counterpart virtual asset service providers
where both sides of a transaction are themselves regulated as such. During the hackathon
phase, testnet transactions do not trigger a live obligation, but the compliance design
should document how this requirement would be satisfied in production.

\subsection*{1.6 Data Privacy}

The retained United Kingdom General Data Protection Regulation and the European Union's
General Data Protection Regulation impose obligations on any entity processing the
personal data of United Kingdom or European Union residents, including a lawful basis for
processing, purpose limitation, data minimisation, a right to erasure, and breach
notification duties \citep{gdpr2016}. The central tension for Blockmediary is that on
chain immutability is structurally incompatible with the right to erasure under Article
17: any personal data written to a public blockchain cannot subsequently be deleted. A
live architectural risk sits within this tension. The full product design commits a
cryptographic hash of the canonical escrow specification to the Base network at the point
a deal is created, and that specification contains party names and know your customer
reference identifiers. Even though a hash is not the underlying plaintext, where
Blockmediary controls both the off chain store containing the names and the on chain hash,
the hash can be linked back to an identified individual, which brings it within the
functional definition of personal data under the General Data Protection Regulation, and
committing that hash to an immutable public ledger creates a permanent record that cannot
be erased, in direct tension with Article 17. This question has not yet been closed and
requires a design decision before the product moves to mainnet. The confirmed
architectural position, however, is that all personal data is kept off chain: the smart
contract stores only wallet addresses, deal amounts, state transitions, and document
hashes, never names, documents, or know your customer records. On chain wallet addresses
are linked to off chain identity records through an internal pseudonymous reference rather
than directly by name. Blockmediary intends to resolve the specification hash problem
before mainnet deployment by stripping personal data from the specification entirely
before it is hashed, which is the simplest of the available options, alongside publishing a
privacy notice, appointing a data protection officer for production, and implementing a
seventy two hour breach notification procedure to the Information Commissioner's Office
consistent with the retained United Kingdom regime.

\subsection*{1.7 Smart Contract Legal Enforceability}

The Law Commission published a report in 2021 confirming that smart contracts can
constitute legally binding agreements under English law, while identifying continuing
uncertainty around offer and acceptance, consideration, mistake, and available remedies
\citep{lawcommission2021}. Blockmediary's smart contract automatically enforces financial
obligations, including fund release and refund, which raises the question of whether the
smart contract itself is the binding agreement between the parties or merely the execution
mechanism for a separate agreement negotiated off chain. A related complication is that on
chain state transitions, once executed, may not be reversible even where a court later
orders that they should be. Blockmediary's response is to maintain a separate signed Trade
Escrow Agreement in natural language that governs the substantive rights of the parties,
treating the smart contract as the automated execution layer for that agreement rather than
the agreement itself. The Trade Escrow Agreement specifies English law as the default
governing law, reflecting its maturity for commercial and trade finance disputes, and
includes an explicit clause addressing the relationship between on chain outcomes and the
underlying agreement, together with a dispute escalation path to an external forum such as
the International Chamber of Commerce's DOCDEX procedure or arbitration under the London
Court of International Arbitration.

\subsection*{1.8 The Electronic Trade Documents Act 2023}

The Electronic Trade Documents Act 2023 came into force in September 2023 and was the
first United Kingdom statute to give electronic trade documents, including electronic
bills of lading, ship's delivery orders, sea waybills, and warehouse receipts, the same
legal status as paper originals under English law \citep{etda2023}. The Act is based on
the United Nations Commission on International Trade Law's Model Law on Electronic
Transferable Records, and equivalent legislation has been adopted in Singapore, Bahrain,
and within the Dubai International Financial Centre, all of which are relevant to
Blockmediary's target corridors \citep{uncitralmletr}. Blockmediary accepts uploaded trade
documents as the basis for releasing escrowed funds, and the legal validity of an
electronic bill of lading in particular, specifically whether it functions as a document
of title giving its holder the right to claim the goods, depends on whether the Act's
requirements are satisfied. Central to this is the Act's control requirement: a person has
control of an electronic trade document only where they can use it, transfer it, and
prevent others from doing so simultaneously, meaning a document that can be freely copied,
such as a portable document format file circulated by email, does not qualify, because it
does not enforce singularity of control. Releasing funds against a document that fails
this test exposes Blockmediary to liability for treating a legally invalid document as
equivalent to an original. For the minimum viable product, Blockmediary's mitigation is to
accept scanned paper documents only, which sidesteps the Act's requirements entirely since
scanned originals remain valid as paper documents, and to defer support for genuine
electronic bills of lading to a post launch phase in which the platform integrates with an
Act qualifying electronic system such as Bolero, essDOCS, or WaveBL, rather than accepting
an ordinary electronic file as if it were equivalent.

\subsection*{1.9 The Financial Conduct Authority's Consumer Duty}

The Consumer Duty, introduced through policy statement PS22/9 and in force since July
2023, substantially expanded the Financial Conduct Authority's previous Principles for
Businesses framework for firms whose products reach retail customers, requiring firms to
deliver good outcomes across four pillars: products and services, price and value,
consumer understanding, and consumer support \citep{fcaconsumerduty}. The Duty applies not
only to firms with a direct retail relationship but to firms anywhere in the distribution
chain through which a retail product reaches an end user. The central question for
Blockmediary is whether its small and medium sized enterprise buyers and sellers count as
retail customers under the Financial Conduct Authority's definition, which excludes only
large institutional counterparties and professional clients meeting specific financial
sophistication thresholds that most small and medium sized enterprises will not meet. If,
as is likely, most of Blockmediary's users are retail customers for these purposes, the
Consumer Duty applies in full, and its obligations go considerably beyond the older
Principles for Businesses framework: the consumer understanding pillar requires that all
communications, including marketing materials, terms of service, and the Trade Escrow
Agreement itself, be tested against whether a reasonable small or medium sized enterprise
in the target market would genuinely understand what it is agreeing to and what risks it is
accepting, and the Duty imposes a proactive obligation to pursue good outcomes rather than
merely to avoid harm. Blockmediary's mitigation is to conduct a formal retail and
professional client classification exercise for its target users, to design the product
and its communications around the consumer understanding standard from the outset, and to
build a complaints handling process that satisfies the Financial Conduct Authority's
dispute resolution rules, including an eight week written response obligation and
referral rights to the Financial Ombudsman Service for eligible complainants. This
obligation applies to the production product; it is not triggered by a testnet build with
no real users.

%---------------------------------------------------------------
\section*{2. European Union Regulatory Risks}
%---------------------------------------------------------------

\subsection*{2.1 The Markets in Crypto Assets Regulation}

The Markets in Crypto Assets Regulation is the European Union's comprehensive framework
for crypto asset activity \citep{mica2023}. Title III governs asset referenced tokens,
stablecoins pegged to multiple assets, requiring authorisation by a national competent
authority, reserve requirements, and whitepaper disclosure. Title IV governs e money
tokens, stablecoins pegged to a single fiat currency, which must be issued by an
authorised credit institution or electronic money institution. Title V governs crypto
asset service providers and requires authorisation to provide services including custody,
exchange, and transfer of crypto assets within the European Union, with full application
from December 2024 and national transitional arrangements expiring on 1 July 2026. USDC is
an e money token under this classification, and Circle, its issuer, must itself hold
Markets in Crypto Assets authorisation to issue USDC within the European Union, which it
obtained through electronic money institution authorisation in France in 2024
\citep{ebamica}. Separately, Blockmediary itself may qualify as a crypto asset service
provider if it provides crypto asset transfer or custody services to European Union
persons, which would require its own Markets in Crypto Assets authorisation, since the
national transitional arrangements that some firms currently rely upon expire on 1 July
2026, after which operating without authorisation for European Union users becomes
unlawful. Blockmediary's mitigation is to confirm with Circle that the specific USDC
deployment it uses remains compliant with the regulation, to apply for crypto asset
service provider authorisation with a national competent authority ahead of any European
Union market entry, which requires an entity established within the European Union, and in
the interim to treat the Abu Dhabi Global Market or the Dubai International Financial
Centre as the initial regulatory home from which European Union authorisation can later be
sought.

\subsection*{2.2 The Digital Operational Resilience Act}

The Digital Operational Resilience Act came into full force on 17 January 2025 and applies
to financial entities operating in the European Union, including crypto asset service
providers authorised under the Markets in Crypto Assets Regulation \citep{dora2022}. The
Act establishes a harmonised framework for information and communication technology risk
management, comprising a mandatory risk management framework, mandatory incident
reporting to competent authorities, digital operational resilience testing including
threat led penetration testing for significant firms, and contractual requirements
governing third party technology providers. The Act is not a standalone licence; rather,
it is a compliance obligation that automatically attaches once a firm holds a Markets in
Crypto Assets authorisation, so the two cannot be addressed separately. Blockmediary's
target architecture depends on several third party technology providers, including
Anthropic for artificial intelligence based document analysis, a roadmap capability not
yet delivered in the current prototype, which instead uses a deterministic rules engine,
cloud infrastructure for document storage and the audit ledger, a know your business and
know your customer provider, and a blockchain node provider for mainnet access, each of
which counts as a third party
information and communication technology provider under the Act, requiring a documented
risk assessment, contractual terms covering audit rights and exit arrangements, and
inclusion in Blockmediary's technology risk register. A significant incident, such as
compromise of the releaser key discussed in section five, a prolonged outage of the
document verification service, or a breach affecting the know your customer database,
would all be reportable under the Act's escalating timeline: an initial notification
within four hours of classification, an intermediate report within seventy two hours, and
a final report within one month. Blockmediary's mitigation is to build its technology risk
framework alongside, rather than after, its Markets in Crypto Assets authorisation
application, to maintain a register of every third party technology provider with its
associated risk classification and contractual terms, and to designate an incident
classification and reporting function, which at minimum viable product stage can sit with
the Chief Compliance Officer.

\subsection*{2.3 Electronic Identification and Trust Services}

The eIDAS 2.0 Regulation, published in April 2024, governs electronic identification and
trust services across the European Union and introduces the European Union Digital
Identity Wallet alongside an updated legal framework for electronic signatures, under
which a qualified electronic signature carries the same legal effect as a handwritten
signature in every member state \citep{eidas2024}. In the United Kingdom, the equivalent
framework is the Electronic Communications Act 2000 \citep{ukeca2000}. The Trade Escrow
Agreement is a binding contract between buyer, seller, and Blockmediary, and its
enforceability across European Union jurisdictions depends on the electronic signature
method used to execute it. A simple click through acceptance or an unconfirmed email may
not satisfy the higher signature standard some member states require for certain contract
types, and a wallet signature produced through signing in with an Ethereum address does
not currently qualify as a qualified electronic signature because it is not issued by a
qualified trust service provider, although it is arguably an advanced electronic signature
in the sense that it is uniquely linked to the signatory's private key and can detect
subsequent tampering. Blockmediary's mitigation is to specify explicitly in the Trade
Escrow Agreement which signature method has been used for a given deal, whether click
through acceptance, wallet signature, or a third party signing platform, and to offer an
advanced electronic signature option through a qualified trust service provider such as
DocuSign or Signicat for higher value or contested deals. Under the United Kingdom's more
functional approach in the Electronic Communications Act 2000, both click through
acceptance and wallet signatures should be treated as valid provided the parties clearly
intend to be bound by them, and the Trade Escrow Agreement includes an explicit clause
confirming that intention.

%---------------------------------------------------------------
\section*{3. Turkey: An Excluded Jurisdiction}
%---------------------------------------------------------------

Turkey is excluded from Blockmediary's target markets, and this section records why
rather than analysing the jurisdiction as an active target. Three instruments drive the
exclusion. The Central Bank of the Republic of Turkey's Regulation No. 2021/14 bans the
use of crypto assets, including stablecoins, as a means of payment for goods and services
within Turkey \citep{cbrt202114}. This is the primary hindrance, because it prohibits
precisely what Blockmediary's product does: using a stablecoin to settle a commercial
transaction. Stablecoins may still be held and traded on licensed platforms in Turkey,
but they cannot be used to pay for a shipment, which makes the core product unusable for
any deal in which a Turkish resident or entity is a buyer or seller. Second, Law No. 7518,
which amended Turkey's Capital Markets Law in 2024, brings crypto asset service providers
under the supervision of the Capital Markets Board and sets a minimum capital requirement
of one hundred million Turkish lira, equivalent to roughly two point three million pounds
sterling, for platforms seeking authorisation, with a certificate of authorisation
required by 30 June 2026 \citep{cmblaw7518}. This capital threshold is a further hindrance
independent of the payments ban, since it would require a dedicated Turkish entity holding
substantial capital purely to serve a market Blockmediary cannot legally transact in under
the payments ban regardless. Third, Turkey's Personal Data Protection Law, known as
KVKK, imposes its own cross border data transfer regime that would apply to any know your
customer data collected from Turkish users \citep{kvkk6698}, but this consideration is
secondary and dormant, since the payments ban already excludes Turkey before any data
protection question arises. Taken together, the payments ban makes the product illegal to
operate in Turkey regardless of licensing, the capital requirement makes licensing
prohibitively expensive even if the ban did not exist, and the data protection regime adds
a further compliance layer on top of both. Turkey is therefore excluded from
Blockmediary's target markets as a considered decision rather than an oversight, to be
revisited only if the underlying regulatory position changes materially.

%---------------------------------------------------------------
\section*{4. United Arab Emirates Regulatory Risks}
%---------------------------------------------------------------

\subsection*{4.1 The Mainland Regime and the CBUAE Overlay}

The Central Bank of the United Arab Emirates regulates payment tokens, defined as crypto
assets fully backed by one or more fiat currencies and used for settlement or transfer,
under its Payment Token Services Regulation, effective since August 2024
\citep{cbuaeptsr2024}. Any entity issuing, redeeming, or facilitating payment tokens on
the mainland must hold a Central Bank licence, and Federal Decree Law No. 6 of 2025
extends this oversight to decentralised finance, cross chain bridges, and broader Web3
infrastructure, with a one year transition period running to September 2026
\citep{uaefdl62025}. Facilitating USDC denominated escrow for mainland parties may
therefore require a Central Bank licence, such as a stored value facility or payment
services licence, and administrative penalties for unlicensed activity reach one billion
United Arab Emirates dirhams. Blockmediary's mitigation is to avoid the mainland regime
entirely at this stage by routing launch activity through the Dubai International
Financial Centre and a separate Virtual Assets Regulatory Authority licensing track,
discussed in sections 4.2 and 4.3.

This mainland avoidance does not, however, close the question entirely. Article 2 of the
Payment Token Services Regulation excludes Financial Free Zones from its definition of the
United Arab Emirates, but a locally licensed virtual asset service provider may still need
to apply for a non-objection registration to custody or transfer certain foreign payment
tokens \citep{cbuaeptsrarticle2}. Whether a United Arab Emirates buyer may lawfully prefund
USDC or EURC into Blockmediary's escrow to settle a physical goods purchase, and whether
the Dubai International Financial Centre partner or the Virtual Assets Regulatory
Authority entity separately needs a Central Bank licence, registration or non-objection
for that specific flow, remains an open legal gate that differs depending on whether the
transaction is structured within the Dubai International Financial Centre, on the
mainland, or cross border to a United Arab Emirates person. Written advice on this point
should be obtained before any live United Arab Emirates launch \citep{difcvarareport2026}.

\subsection*{4.2 The Dubai International Financial Centre Partner-Led Pilot}

Following further discussion within the team, including an existing relationship with a
Dubai International Financial Centre connected partner, Blockmediary's near term United
Arab Emirates route is a partner-led pilot within the Centre rather than an independent
licence application. A Dubai International Financial Centre partner-led pilot is not
itself a named Dubai Financial Services Authority licence category. It is an operating
arrangement in which a Dubai Financial Services Authority authorised firm acts as the
regulated principal and Blockmediary supplies the technology, document verification and
workflow layer beneath it. A Dubai International Financial Centre Commercial or Innovation
Licence gives Blockmediary a legal presence in the Centre but does not itself authorise
regulated financial services \citep{difcvarareport2026}. Under the Authority's outsourcing
rules, the authorised partner remains responsible for the outsourced functions, must
conduct due diligence on and supervise Blockmediary as provider, and, where the
arrangement is material, must notify the Authority and maintain written outsourcing and
contingency arrangements \citep{dfsaoutsourcing}. The partner agreement and responsibility
matrix are therefore the central regulatory instrument for the pilot, not a licence held
by Blockmediary itself.

Minimum requirements before the pilot may accept live funds include a written regulatory
perimeter opinion mapping the customer journey, entities, wallet flows, release key and
smart contract control; verification that the partner's actual Dubai Financial Services
Authority permissions cover the specific custody, settlement and customer activity
involved, not merely generic fintech or payment permissions; an executed partner term
sheet and outsourcing agreement covering scope, audit rights, incident reporting and
termination; a written responsibility matrix assigning customer contracting, onboarding,
sanctions screening, wallet screening, safeguarding, release approval and loss liability;
and a defined pilot envelope limited to business to business customers, deal values at or
below the fifty thousand pound automation threshold, approved low risk corridors, and
manual sign off available on every release. The Dubai Financial Services Authority's
Innovation Testing Licence is not assumed to be required or sufficient for this route:
it is a restricted sandbox for a firm that itself needs to test a regulated activity, and
should be pursued only if the Authority or counsel concludes Blockmediary will itself
conduct regulated activity during testing, not treated as a default or cheap substitute
for a properly structured partner arrangement \citep{difcvarareport2026}.

\subsection*{4.3 Virtual Assets Regulatory Authority Licensing for Scale}

Alongside, and starting from the same point in time as, the Dubai International Financial
Centre pilot, Blockmediary pursues its own Virtual Assets Regulatory Authority licence for
wider Dubai scale. The Authority regulates virtual asset activity outside the Dubai
International Financial Centre, and an application must be made through a non-Centre
entity, meaning Blockmediary needs a separate legal entity or group company established
outside the Centre, with a physical Dubai office and two full time, fit and proper
Responsible Individuals \citep{difcvarareport2026}. The Authority licenses each activity
separately, and the closest apparent fit for Blockmediary is Virtual Asset Transfer and
Settlement, covering the transmission or settlement of virtual assets from one entity or
wallet to another, which matches the product's release of escrowed stablecoin to the
seller once documentary conditions are satisfied. Whether Custody Services also applies is
conditional rather than settled: direct smart contract funding helps, but the server side
release key and any administration, upgrade or recovery rights Blockmediary retains may
still amount to control or safekeeping requiring a written opinion. Broker Dealer
Services, covering arranging or matching orders to buy or sell virtual assets, is not
established by the current model, since Blockmediary settles a physical goods sale rather
than arranging a token purchase or sale, and should not be assumed mandatory
\citep{difcvarareport2026}. If Custody is ultimately required, the Authority's rules
generally require Custody to sit in a legally separate entity from other virtual asset
activities, with only limited scope to combine Custody and Transfer and Settlement subject
to the Authority's approval and strict operational segregation, so the assumption that one
entity will hold all three activities should not be treated as settled.

The Authority does not publish a guaranteed approval timeline. The two workstreams start
together at inception, modelled as Month 1: the Dubai International Financial Centre pilot
targets its first paid transaction in Month 12, and the Virtual Assets Regulatory Authority
application targets a full licence in Month 18 in the base case, with Month 12 retained
only as an early case where the activity perimeter resolves quickly and remediation is
minimal, and Month 24 as a delayed case reflecting regulatory or partner remediation.
Approval to Incorporate or in principle approval at an earlier stage is not permission to
serve customers, and Virtual Assets Regulatory Authority scale activity should begin only
from the effective date of the full licence \citep{difcvarareport2026}. Indicative fees
for a broad, three activity scope sit around one hundred and seventy thousand dirhams in
application fees, four hundred and eighty thousand dirhams a year in supervision, and a
capital floor of approximately one and a half million dirhams, though a narrower activity
scope, for instance Transfer and Settlement alone, would reduce each figure materially; a
Transfer and Settlement only scope is closer to forty thousand dirhams in application
fees, eighty thousand dirhams a year in supervision, and a capital floor of the higher of
five hundred thousand dirhams or twenty five per cent of fixed annual overhead
\citep{vararulebook2025fees}. These remain conservative planning placeholders pending a
written activity opinion, not confirmed fees.

\subsection*{4.4 Data Protection}

The United Arab Emirates' federal Personal Data Protection Law applies to any entity
processing the personal data of United Arab Emirates residents regardless of where that
entity is located, and requires a lawful basis for processing, data minimisation, a data
protection officer for high volume sensitive processing, and breach notification
\citep{pdpl2021}. Executive regulations remain pending as of early 2026, with a six month
implementation grace period expected once they are published. Know your business and know
your customer data collected from United Arab Emirates users, including passports,
business registration documents, and wallet data linked to identity, is personal data
under this law. The Dubai International Financial Centre maintains its own Data
Protection Law of 2020, arguably more mature than the federal law and the relevant
framework for the Centre entity conducting the partner-led pilot, and it requires its own
notification where applicable \citep{difcdplaw2020}. The separate Virtual Assets
Regulatory Authority entity, established outside the Centre, falls instead under the
federal law directly. Blockmediary's mitigation is to align with the federal law's core
principles across both entities now, to comply with the Centre's own framework and
notification requirements for the Dubai International Financial Centre entity
specifically, and to monitor the federal executive regulations once published so that any
additional requirements can be implemented within the grace period.

%---------------------------------------------------------------
\section*{5. Cross Cutting Risks}
%---------------------------------------------------------------

\subsection*{5.1 Sanctions Screening}

Because Blockmediary's model spans multiple trade corridors, every deal potentially
involves a party, a wallet, a category of goods, or a corridor subject to sanctions
imposed by the Office of Foreign Assets Control, the United Nations Security Council, His
Majesty's Treasury, or the United Arab Emirates. Blockmediary's mitigation is to screen
both parties, all wallets, the category of goods, and the trade corridor against every
applicable sanctions list before allowing escrow funding, to re-screen at the point of
document review and again at the point of fund release, since sanctions can be imposed
between deal creation and settlement, to maintain a blocked countries, blocked goods, and
blocked corridors list, and to escalate any hit to the Chief Compliance Officer before
proceeding under any circumstances.

\subsection*{5.2 Trade Based Money Laundering}

Cross border trade escrow is a recognised vector for trade based money laundering. The
established red flags include an invoice value inconsistent with the market price of the
goods described, a vague goods description such as general merchandise, unusual trade
routes or transshipment points, shipment to or from a high risk jurisdiction, repeated
transactions between the same parties with changing terms, and alterations or
inconsistencies across successive versions of a submitted document. Blockmediary's
mitigation is to build a trade based money laundering red flag checklist into the document
review workflow, to add artificial intelligence assisted document extraction once that
roadmap capability is delivered to flag price and quantity inconsistencies automatically,
to require mandatory human review before any release involving a first time counterparty
pairing, and to train document reviewers on the Financial Action Task Force's published
typologies for this form of laundering. The current prototype's deterministic rules engine
already performs the equivalent checks on submitted structured fields.

\subsection*{5.3 Governing Law, Jurisdiction, and Dispute Forum}

Without a clear governing law and dispute resolution clause, a dispute between a buyer in
one jurisdiction and a seller in another over a Blockmediary structured escrow risks
jurisdictional confusion. Blockmediary's mitigation is to default to English law as the
governing law, given its maturity for commercial and trade finance disputes, to default to
the International Chamber of Commerce's DOCDEX procedure for documentary disputes and to
London Court of International Arbitration or International Chamber of Commerce
arbitration for larger commercial disputes \citep{iccdocdex}, while allowing parties to
elect an alternative governing law or forum within the set of jurisdictions Blockmediary
supports, with the chosen clause recorded explicitly in the Trade Escrow Agreement
template.

\subsection*{5.4 Smart Contract Security and Audit}

A bug or exploit in the escrow smart contract could result in the loss of user funds, and
regulators including the Financial Conduct Authority, the Virtual Assets Regulatory
Authority, and the Turkish Capital Markets Board increasingly expect cryptoasset firms to
demonstrate security standards. During the hackathon, the mitigation is limited to code
review and peer testing, with the build clearly marked as a testnet minimum viable
product, but production deployment requires an independent smart contract audit from a
reputable firm before real funds are placed at risk, alongside role based access control,
a multi signature requirement for the release authority, and a time lock mechanism, backed
by a bug bounty programme once the audit is complete.

Within this category, releaser key compromise is the single most severe technical risk
Blockmediary faces. The releaser role's private key is the sole bridge between an off
chain compliance verdict and the corresponding on chain state transition, meaning that an
attacker who compromises this key can move any deal to a release pending state regardless
of whether the underlying documents were actually compliant, can redirect funded deals
toward a refund, and cannot be stopped once a malicious transaction has been submitted to
the chain. This is not merely a technical incident; it is a customer funds loss event that
would trigger breach notification to the Financial Conduct Authority under the payment
services or cryptoasset regime, incident reporting to the Virtual Assets Regulatory
Authority within twenty four hours under its Rulebook, potential criminal liability under
the Proceeds of Crime Act, and civil liability to the affected buyers and sellers.
Blockmediary's mitigation for production is to hold the releaser key in a hardware
security module or a threshold signature scheme rather than a conventional hot wallet, to
implement role rotation so that the releaser key can be replaced immediately on suspected
compromise without redeploying the contract, to pause the contract immediately on
suspected compromise using a separate, ideally multi signature, administrator key, and to
maintain a documented incident response sequence running from pause, to revocation, to
rotation, to an audit ledger review, and finally to regulatory and user notification.

A related risk is that the same pause capability that protects against key compromise can
also function as a censorship lever if misused: once a deal has reached a release pending
state, meaning the documents have been verified as compliant and any objection window has
expired without a valid objection, the seller has a contractual right to payment, and
invoking a pause at that point exposes Blockmediary to a breach of contract claim.
Blockmediary's mitigation is to define exhaustive, narrow grounds for pausing in the Trade
Escrow Agreement, limited to a suspected exploit or critical vulnerability, a regulatory
direction or court order, or a confirmed sanctions hit discovered after the verdict was
recorded, to treat any other use of the pause function as a breach, to require joint
approval from the Chief Compliance Officer and the Chief Technology Officer before a pause
is invoked, and to notify both parties immediately with the stated reason and an estimated
resolution timeline whenever a pause does occur.

\subsection*{5.5 Documentary Credit Standards}

Blockmediary positions its product as applying the logic of a documentary credit without
requiring an issuing bank, so the International Chamber of Commerce's Uniform Customs and
Practice for Documentary Credits, universally known as UCP 600, functions as an inherited
operational standard even though Blockmediary's Trade Escrow Agreement is a distinct
instrument and not itself a letter of credit \citep{ucp600}. A number of UCP 600 articles
translate directly into requirements for Blockmediary's document verification engine.
Articles 4 and 5 establish the autonomy principle, under which a documentary credit deals
only in documents rather than in the underlying goods or performance, meaning that once
Blockmediary's rules engine confirms compliant documents, funds must release even where
the buyer separately claims the goods themselves are defective, a position that must be
disclosed clearly in the Trade Escrow Agreement. Article 14(a) sets a five banking day
examination window, which Blockmediary must replicate as a hard review deadline, since
missing it precludes the platform from later treating documents as non compliant. Article
14(b) permits data to differ in wording across documents provided it does not conflict,
so Blockmediary's automated checker should flag genuine conflicts, such as a quantity
stated as five thousand metric tonnes on the invoice against five metric tonnes on the
bill of lading, rather than immaterial differences in phrasing. Article 14(c) requires a
transport document to be presented within twenty one days of the shipment date, which
Blockmediary must validate against the stated shipment date. Article 14(d) requires the
goods description on the invoice to match the underlying agreement exactly, while other
documents need only avoid inconsistency, a distinction the verification engine must
preserve rather than applying a uniformly strict or uniformly loose standard. Article
14(g) provides that a condition without a named supporting document is disregarded, so
every compliance condition Blockmediary builds into a deal must map to a specific
document. Article 15 makes a complying presentation mandatory to honour, removing any
discretion to delay release once compliance is confirmed, and Article 16 requires that
refusal be accompanied by an itemised discrepancy notice within the same five banking day
window, with late notice forfeiting the right to rely on non compliance, which
Blockmediary's rejection workflow must mirror, including a mechanism for the buyer to
waive a discrepancy. Article 17 requires at least one original of each document, which the
upload flow must enforce, with electronic documents governed separately by the eUCP
supplement \citep{eucp}. Article 18 sets four checks for a commercial invoice: that it is
issued by the seller, that the buyer is the named consignee, that the currency matches the
deal currency, a point of particular importance for a USDC denominated escrow, and that
the goods description matches exactly. Article 20 sets the corresponding checks for a bill
of lading: an on board notation with a date, a named vessel, a clean document with no
damage clause, and matching ports of loading and discharge, all of which the verification
engine must parse. Article 27 makes clear that any damage or defect clause on a transport
document removes it from automatic processing and requires human review. Article 28
requires that where insurance is required, typically under a cost, insurance, and freight
or carriage and insurance paid term, cover must be at least one hundred and ten per cent
of the relevant value, must extend to the destination, and cannot be satisfied by a mere
cover note. Article 30 sets tolerances of plus or minus ten per cent where a quantity is
qualified as about or approximately, and plus or minus five per cent otherwise, which the
verification engine's tolerance logic must implement directly.

The International Chamber of Commerce's Incoterms 2020 rules operate alongside UCP 600 and
determine which documents a given deal requires, since the Incoterm agreed between buyer
and seller fixes where risk transfers, who bears freight and insurance costs, and which
documents the seller must produce \citep{incoterms2020}. Getting the document checklist
wrong for the stated Incoterm either exposes the buyer to an incomplete presentation being
accepted or exposes the seller to a compliant presentation being wrongly rejected. Terms
requiring insurance, cost, insurance, and freight and carriage and insurance paid, demand
an insurance certificate rather than the cover note UCP 600 Article 28(c) prohibits, and
carriage and insurance paid specifically requires cover to Institute Cargo Clauses A, the
all risks standard the 2020 revision introduced in place of the narrower C clauses used
previously. Delivered duty paid terms require evidence of import customs clearance, the
absence of which produces an incomplete presentation. Free on board, cost and freight, and
cost, insurance, and freight all require an on board bill of lading, and although these
sea and inland waterway terms are technically inappropriate for containerised cargo, they
remain widely used in practice, so Blockmediary's verification engine must accommodate the
practice while flagging the technical mismatch for review. The 2020 revision also renamed
delivered at terminal to delivered at place unloaded and permits the buyer's own carrier to
issue an on board bill of lading to the seller under free carrier terms, resolving a long
standing mismatch between that term and documentary credit practice, and Blockmediary
should build this election into deal setup as an option. Ex works terms are not recommended
on the platform for cross border deals, because the seller bears almost no documentary
obligation under that term, which makes Blockmediary's release logic largely unworkable.
Blockmediary's mitigation across both standards is to require the Incoterm to be specified
at deal setup and to generate the document checklist dynamically from it, to validate that
the transport mode matches the Incoterm selected, to check that any required insurance
certificate meets the correct Institute Cargo Clauses standard and value threshold, and to
include a representation in the Trade Escrow Agreement that the stated Incoterm governs
the parties' documentary obligations without any duty on Blockmediary's part to verify
physical delivery beyond the documentary evidence presented.

\subsection*{5.6 Export Controls}

Independently of sanctions, several export control regimes impose a positive obligation to
obtain a licence before shipping certain categories of goods, regardless of whether either
counterparty is sanctioned. The United Kingdom's Export Control Order 2008 covers military
goods and dual use goods mirroring the European Union's dual use list, administered by the
Export Control Joint Unit \citep{ukeco2008}. The European Union's Dual Use Regulation,
updated in 2021, covers goods, software, and technology with both civilian and military
applications, including encryption, certain chemicals, optical equipment, and electronics
\citep{eudualuse2021}. The United States Export Administration Regulations and
International Traffic in Arms Regulations apply to United States origin goods, United
States technology, and transactions involving United States persons regardless of where
Blockmediary itself is incorporated, meaning that a seller shipping United States origin
goods subject to the Export Administration Regulations may need a licence from the Bureau
of Industry and Security even where the transaction is structured through a United Kingdom
or United Arab Emirates entity \citep{usear, usitar}. A cross border escrow platform that
releases payment for a shipment of export controlled goods without checking for a valid
export licence risks facilitating an illegal export, a criminal offence under each of
these regimes. Dual use goods are not confined to obviously military items; they include
encryption software, certain chemicals, precision optics, drone components, thermal
imaging equipment, and high performance computers, categories a small or medium sized
enterprise exporter may not recognise as controlled. Blockmediary's target market already
excludes sanctioned corridors and high risk regulated goods in principle, but this
exclusion requires an operational mechanism rather than a policy statement alone.
Blockmediary's mitigation is to add an export control goods category screen to the deal
intake compliance gate, distinct from sanctions screening, collecting the goods
description, the Harmonised System code, and the origin and destination countries at
intake, and cross referencing these against the United Kingdom and European Union dual use
lists and the United States Commerce Control List, requiring the seller to confirm and
evidence a valid export licence where the goods fall on a controlled list, maintaining a
controlled goods category list within the platform's rules engine, and consulting export
control counsel before onboarding any corridor involving United States origin goods or
United States persons.

%---------------------------------------------------------------
\section*{6. Regulatory Reference Index}
%---------------------------------------------------------------

The instruments referenced throughout this assessment are consolidated here by
jurisdiction. In the United Kingdom, the relevant instruments are the Financial Conduct
Authority's consultation papers CP25/14 on stablecoin issuance and cryptoasset custody and
CP25/40 on the broader regulation of cryptoasset activities \citep{fcacp2514, fcacp2540},
the Payment Services Regulations 2017 \citep{ukpsr2017}, the Money Laundering Regulations
2017 as amended \citep{mlr2017}, the retained United Kingdom General Data Protection
Regulation \citep{gdpr2016}, the Law Commission's 2021 report on smart legal contracts
\citep{lawcommission2021}, the Electronic Trade Documents Act 2023 \citep{etda2023}, the
Consumer Duty under policy statement PS22/9 \citep{fcaconsumerduty}, the Electronic
Communications Act 2000 \citep{ukeca2000}, the Proceeds of Crime Act 2002
\citep{poca2002}, the Terrorism Act 2000 \citep{terrorismact2000}, and the Export Control
Order 2008 \citep{ukeco2008}.

In the European Union, the relevant instruments are the Markets in Crypto Assets
Regulation \citep{mica2023} and the European Banking Authority's guidance on asset
referenced and e money tokens issued under it \citep{ebamica}, the General Data Protection
Regulation \citep{gdpr2016}, the sixth Anti Money Laundering Directive \citep{amld6}, the
Digital Operational Resilience Act \citep{dora2022}, the eIDAS 2.0 Regulation
\citep{eidas2024}, and the Dual Use Regulation \citep{eudualuse2021}.

In Turkey, the relevant instruments, retained here for reference despite the
jurisdiction's exclusion, are Law No. 7518 governing capital markets and crypto asset
service provider licensing \citep{cmblaw7518}, Central Bank Regulation No. 2021/14
banning crypto asset payments \citep{cbrt202114}, and the Personal Data Protection Law
known as KVKK \citep{kvkk6698}.

In the United Arab Emirates, the relevant instruments are the Central Bank's Payment
Token Services Regulation and its Article 2 exclusion of Financial Free Zones
\citep{cbuaeptsr2024, cbuaeptsrarticle2}, Federal Decree Law No. 6 of 2025 extending
oversight to decentralised finance and Web3 infrastructure \citep{uaefdl62025}, the Virtual
Assets Regulatory Authority's Rulebook version 2.0, its Schedule 1 activity definitions,
and its Schedule 2 supervision and authorisation fees \citep{vararulebook2025,
vararulebook2025activities, vararulebook2025fees}, the federal Personal Data Protection Law
\citep{pdpl2021}, the Dubai International Financial Centre's Data Protection Law
\citep{difcdplaw2020}, the Dubai Financial Services Authority's crypto token regime and its
outsourcing rules under GEN 5.3.21 to 5.3.22 \citep{difcdfsa2022, dfsaoutsourcing}, and the
internal readiness review reconciling all three regulators against Blockmediary's product
architecture \citep{difcvarareport2026}. The Abu Dhabi Global Market's Data Protection
Regulations and the Financial Services Regulatory Authority's digital asset guidance
\citep{adgmdpr2021, adgmfsra2026} are retained here only because the Abu Dhabi Global
Market was considered and moved away from, as recorded in section 8; they are not part of
the confirmed launch route.

Globally applicable instruments include the Financial Action Task Force's Recommendation
16 on the travel rule \citep{fatf2025}, the Uniform Customs and Practice for Documentary
Credits, UCP 600, and its electronic supplement, the eUCP \citep{ucp600, eucp}, Incoterms
2020 \citep{incoterms2020}, the International Chamber of Commerce's DOCDEX dispute
resolution procedure \citep{iccdocdex}, the United Nations Commission on International
Trade Law's Model Law on Electronic Transferable Records \citep{uncitralmletr}, and the
United States Export Administration Regulations and International Traffic in Arms
Regulations \citep{usear, usitar}.

%---------------------------------------------------------------
\section*{7. Competitive Landscape: Compliance Comparison}
%---------------------------------------------------------------

Understanding how the closest competing platforms handle compliance and release logic
sharpens Blockmediary's own positioning. Tazapay operates one of the largest cross border
payment and escrow infrastructure platforms for business to business trade, collecting
payments from buyers, holding them in escrow, and releasing them to sellers once agreed
conditions are met, and it also supports conversion between stablecoins and fiat currency
\citep{tazapay2025}. It operates across seventy markets for collections and pays out to
over one hundred countries, and is licensed in Singapore, Canada, Australia, and the
United States, with licence applications pending in the United Arab Emirates and the
European Union. Its release model is buyer discretion: although Tazapay checks some
documents, funds are ultimately released when the buyer clicks approve rather than through
an independent documentary compliance check. This is the key distinction from
Blockmediary, whose release is triggered by document compliance regardless of whether the
buyer separately approves, a mechanism much closer to how a letter of credit functions,
and Tazapay has no on chain smart contract component at all.

Truzo is a regulated digital escrow and payment platform originally focused on the United
Kingdom to South Africa trade corridor, and is notable as the first Financial Conduct
Authority approved digital escrow service focused on Africa, partnering with Currencycloud
for cross border currency conversion \citep{truzo2023}. It is smaller and corridor
focused, backed by the United Kingdom's Department for International Trade, and operates
in fiat only. Its release model is traditional escrow, holding funds in a bank account and
releasing them on delivery confirmation, with no blockchain or smart contract component.
Truzo is, in substance, a digital version of a conventional escrow agent running on
ordinary banking rails, whereas Blockmediary targets a broader set of corridors using on
chain stablecoin escrow with documentary release logic.

XREX operates a Taiwanese crypto financial platform whose cross border escrow product,
BitCheck, escrows stablecoins and other cryptocurrencies between buyers and sellers for
business to business transactions, with all users passing through know your customer and
anti money laundering verification \citep{xrex2025}. It is licensed by the Monetary
Authority of Singapore, registered as a money services business in the United States, and
has entered Lithuania for European Union market access, though it holds no confirmed
licence in the United Kingdom or the United Arab Emirates. Its release model, like
Tazapay's, is buyer discretion: the seller requests release after shipping and the buyer
then approves it, meaning the buyer can still delay or block payment even where the
underlying documents are fully compliant. XREX is the closest competitor in the sense that
it also uses stablecoin escrow for cross border trade, but its release model is buyer
approval rather than documentary compliance, which is precisely the gap Blockmediary's
Uniform Customs and Practice for Documentary Credits inspired release logic closes: once
documents are verified compliant, funds must release and the buyer cannot simply refuse,
which is what makes Blockmediary a genuine letter of credit alternative rather than a
digital hold and release service. XREX's primary markets in Taiwan, Singapore, and
Southeast Asia also do not significantly overlap with Blockmediary's target markets in the
United Kingdom, the European Union, and the United Arab Emirates.

Komgo operates a blockchain based platform for commodity trade finance, principally oil,
gas, metals, and agricultural commodities, digitising and accelerating the existing letter
of credit process rather than replacing it, and its principal users are large commodity
traders and banks, backed by institutions including Societe Generale, ING, ABN AMRO, and
Citi \citep{komgo2020}. It runs on JPMorgan's Quorum, an Ethereum variant, connected to the
Vakt post trade platform, and uses smart contracts to automate letter of credit matching
and processing, cutting turnaround from roughly ten days to roughly one hour in production.
Critically, however, the letter of credit itself remains the underlying legal instrument;
Komgo is the digital rail beneath it rather than a replacement for it. Because Komgo
modernises letters of credit for large commodity traders and banks rather than serving
small and medium sized enterprises directly, it is not a direct competitor to
Blockmediary, though its document verification and smart contract architecture remain
worth studying as a reference point and a potential future integration partner for
institutional commodity trade flows.

Contour represented the most ambitious attempt to replace the traditional letter of credit
with a blockchain based alternative, backed by nine major banks including HSBC, Standard
Chartered, Citi, and BNP Paribas, and built on R3's Corda blockchain, before winding down
in November 2023 \citep{contour2023}. The stated reasons for its closure were that its
bank shareholders withdrew funding, the platform was processing only sixty to seventy
transactions per month at the point of closure, no single lead investor existed to drive
strategic direction, and the surrounding regulatory frameworks took too long to catch up
with what the platform was attempting to build. Contour's failure left the letter of
credit replacement space open rather than closed. Its core lesson for Blockmediary is
architectural: Contour tried to bring banks onto a new shared network, whereas Blockmediary
bypasses the banking layer entirely through a smart contract model that does not depend on
bank participation to function. Because pitch audiences and regulators alike are likely to
already know Contour's history, framing Blockmediary as having learned from that specific
failure, rather than simply asserting novelty, is a stronger and more credible narrative.

Across all five comparators, the pattern that emerges is consistent: none of Tazapay,
Truzo, or XREX releases funds on the basis of documentary compliance, all three instead
relying on some form of buyer discretion or delivery confirmation, Komgo digitises the
letter of credit rather than replacing it and is not aimed at small and medium sized
enterprises, and Contour attempted a genuine replacement but failed for governance and
funding reasons rather than technical ones. Blockmediary's defensible differentiator is
therefore documentary release itself: funds move when the submitted documents are verified
as compliant, not when a counterparty elects to approve them, which is the specific
mechanism that closes the gap left open by every one of these five comparators.

%---------------------------------------------------------------
\section*{8. United Arab Emirates Launch Route: DIFC Pilot and VARA Scale}
%---------------------------------------------------------------

This section previously weighed a single choice of regulatory home among the Jebel Ali
Free Zone Authority, the Dubai International Financial Centre, and the Abu Dhabi Global
Market, and recommended incorporating in the Abu Dhabi Global Market first. That
recommendation has since been superseded. Following a team decision informed by an
existing relationship one of the founders holds with a Dubai International Financial
Centre connected partner, and by the more detailed readiness analysis summarised in
sections 4.2 and 4.3, Blockmediary's confirmed route is a Dubai International Financial
Centre partner-led pilot run in parallel with an independent Virtual Assets Regulatory
Authority licence application, rather than a single free zone choice or an Abu Dhabi Global
Market first sequence. The Abu Dhabi Global Market is not part of the confirmed plan and is
not discussed further here; it is recorded as a route the team considered and moved away
from, for transparency, rather than omitted silently.

The two tracks are legally distinct and start together at inception, modelled as Month 1.
The Dubai International Financial Centre track, detailed in section 4.2, gives
Blockmediary an early, controlled route to paying customers under an authorised partner's
permissions, targeting a first paid transaction in Month 12. It does not by itself make
Blockmediary a regulated entity, and Blockmediary should not describe itself as Dubai
Financial Services Authority authorised unless and until it separately obtains that
authorisation. The Virtual Assets Regulatory Authority track, detailed in section 4.3,
builds Blockmediary's own licensed platform for wider Dubai scale through a legal entity
established outside the Centre, targeting a full licence in Month 18 in the base case.
Running both from Month 1 avoids making the scale track wait for pilot revenue, while
retaining the pilot as the nearer term route to real customer evidence.

The Jebel Ali Free Zone Authority remains relevant only in the same narrow role
identified previously: it is a logistics and trading free zone, not a financial services
regulator, and could house a holding or commercial subsidiary for physical goods related
activity if the team wanted one, but it has no role in either the Dubai International
Financial Centre pilot or the Virtual Assets Regulatory Authority licence and is not
discussed further.

On data protection, the confirmed route means the Dubai International Financial Centre's
own Data Protection Law of 2020 is the operative free zone framework, since the Centre
entity is the one conducting the partner-led pilot, while the separate Virtual Assets
Regulatory Authority entity, established outside the Centre, falls under the federal
Personal Data Protection Law directly rather than under a second free zone framework, as
set out in section 4.4. European Union residents' data remains subject to the General
Data Protection Regulation directly regardless of where Blockmediary is incorporated, so
neither entity's location exempts the product from that separate obligation.

Three points of continuing uncertainty, carried over directly from the readiness report,
should be tracked as open items rather than treated as resolved by this route decision.
First, whether Blockmediary's release key, contract administration and recovery design
amount to custody or control under either the Dubai Financial Services Authority's or the
Virtual Assets Regulatory Authority's rules is not yet settled and determines entity
structure and capital. Second, the exact Virtual Assets Regulatory Authority activity
scope, and in particular whether Custody Services applies alongside Transfer and
Settlement, is not yet confirmed and affects fees, capital and whether a further separate
entity is required. Third, the Central Bank of the United Arab Emirates treatment of
United Arab Emirates denominated USDC or EURC settlement remains a hard legal gate
independent of which free zone or licence route is chosen, discussed in section 4.1. The
single most valuable next step identified in the readiness review is a single integrated
perimeter memorandum, prepared with counsel and covering the Dubai Financial Services
Authority, Virtual Assets Regulatory Authority and Central Bank positions together, rather
than three separate opinions \citep{difcvarareport2026}.

\newpage
\bibliographystyle{apacite}
\bibliography{Risk}

\end{document}
